\begin{prob}
    For each of the following pieces of code, state the time complexity using
    $\Theta$ notation in as simple of terms as possible. You do not need to
    show your work (but doing so might help you earn partial credit in case your
    overall answer is not correct).

    \begin{subprobset}

        \begin{subprob}
            \inputminted{python}{\thisdir/include/part_a.py}

            \begin{soln}
                $\Theta(n^2)$.

                This is very similar to the dependent loops we saw in the
                tallest doctor problem where we considered unordered pairs.
            \end{soln}
        \end{subprob}

        \begin{subprob}
            \inputminted{python}{\thisdir/include/part_b.py}

            \begin{soln}
                $\Theta(n^15)$.

                Although we have a \python{while}-loop, it essentially
                behaves like a nested \python{for}-loop over \python{range(n)}.
                So we have independent loops, and the total number of executions
                of the innermost loop body is $\Theta(n^15)$.
            \end{soln}
        \end{subprob}

        \begin{subprob}
            \inputminted{python}{\thisdir/include/part_c.py}

            \begin{soln}
                $\Theta(n^2)$.

                When analyzing the time complexity of a piece of code where
                each line takes constant time to execute once, we just need to find the line
                that executes the most number of times. In this case, it's the
                inner body of the nested \python{for}-loop. It executes $n^2$ times, for
                a time complexity of $\Theta(n^2)$.

                We don't know if the last $\python{for}$ loop will run or not
                -- that depends on whether $n$ is even or odd. But even if it's
                even (no pun intended) and it \emph{does} run, it won't change
                the overall time complexity because $n^2 + n$ is still
                $\Theta(n^2)$.
            \end{soln}
        \end{subprob}

        \begin{subprob}
            \inputminted{python}{\thisdir/include/part_d.py}

            \begin{soln}
                $\Theta(n^2)$.

                This is an example of a piece of code where not every line
                takes constant time to execute! In particular, the \python{sum(arr)} line
                takes $\Theta(n)$ time to execute once, and it is executed $n$ times, for
                a total time of $\Theta(n^2)$.
            \end{soln}
        \end{subprob}

        \begin{subprob}
            \inputminted{python}{\thisdir/include/part_e.py}

            \begin{soln}
                This outer \python{for}-loop makes 4 iterations -- one for each
                element of \python{ranges}.

                On the first iteration of the outer loop, the inner loop makes
                10 iterations, taking a total of 10 units of time. That is, the
                first iteration takes constant time.

                On the second iteration of the outer loop, the inner loop makes
                roughly $n/2$ iterations, taking a total of $\Theta(n/2) =
                \Theta(n)$ time.

                On the third iteration of the outer loop, the inner loop makes
                $n$ iterations, taking $\Theta(n)$ time.

                On the fourth iteration of the outer loop, the inner loop makes
                $n^2$ iterations, taking $\Theta(n^2)$ time.

                In total, the time taken is $\Theta(1) + \Theta(n) + \Theta(n) + \Theta(n^2) = \Theta(n^2)$.
            \end{soln}
        \end{subprob}

        \begin{subprob}
            \inputminted{python}{\thisdir/include/part_f.py}

            \begin{soln}
                $\Theta(n^2)$.

                The innermost loop, \python{for i in range(rng_a + rng_b)}, is executed 
                16 times (once for each ordered pair of ranges from \python{ranges}), making
                a different number of executions each time.

                For example, one execution of the innermost loop will be when \python{rng_a = n} and \python{rng_b = n**2}.
                In this case, the innermost loop will make $n + n^2$ iterations, taking $\Theta(n^2)$ time.

                From there, you might see that the overall time complexity will
                be $\Theta(n^2)$, because of the 16 executions of the innermost
                loop, the most costly will be the one where both \python{rng_a}
                and \python{rng_b} are $n^2$. In this situation, the innermost
                loop will make $n^2 + n^2 = 2n^2$ iterations, taking
                $\Theta(n^2)$ time. All of the other executions are
                $\Theta(n^2)$ or smaller, so the total time taken will be
                $\Theta(n^2)$.

                Alternatively, you could write out the time taken by each of the 16 executions of the innermost loop.
                Those 16 executions are:

                % ranges = [10, n/2, n, n**2]

                \begin{center}
                    \begin{tabular}{c|c|c|c}\toprule
                        \python{rng_a} & \python{rng_b} & \# iterations & time taken \\ \midrule
                        10 & 10 & 20 & $\Theta(1)$ \\
                        10 & $n/2$ & $10 + n/2$ & $\Theta(n)$ \\
                        10 & $n$ & $10 + n$ & $\Theta(n)$ \\
                        10 & $n^2$ & $10 + n^2$ & $\Theta(n^2)$ \\
                        $n/2$ & 10 & $n/2 + 10$ & $\Theta(n)$ \\
                        $n/2$ & $n/2$ & $n/2 + n/2$ & $\Theta(n)$ \\
                        $n/2$ & $n$ & $n/2 + n$ & $\Theta(n)$ \\
                        $n/2$ & $n^2$ & $n/2 + n^2$ & $\Theta(n^2)$ \\
                        $n$ & 10 & $n + 10$ & $\Theta(n)$ \\
                        $n$ & $n/2$ & $n + n/2$ & $\Theta(n)$ \\
                        $n$ & $n$ & $n + n$ & $\Theta(n)$ \\
                        $n$ & $n^2$ & $n + n^2$ & $\Theta(n^2)$ \\
                        $n^2$ & 10 & $n^2 + 10$ & $\Theta(n^2)$ \\
                        $n^2$ & $n/2$ & $n^2 + n/2$ & $\Theta(n^2)$ \\
                        $n^2$ & $n$ & $n^2 + n$ & $\Theta(n^2)$ \\
                        $n^2$ & $n^2$ & $n^2 + n^2$ & $\Theta(n^2)$ \\ \bottomrule
                    \end{tabular}
                \end{center}

                Adding up the 16 elements in the rightmost column, we get $\Theta(n^2)$.



            \end{soln}
        \end{subprob}

    \end{subprobset}
    
\end{prob}
